\documentclass{exam}

\newif\ifA
\newif\ifkey
\newif\ifprintselected   % required by exam_config's \ansbox

%%%%%%%%%%%%%%%%%%%%%%%%%%%%%%%%%%%%%%%%%%%%%%%%%%%%%%
%%% SET THESE IF VARIABLES
%%%%%%%%%%%%%%%%%%%%%%%%%%%%%%%%%%%%%%%%%%%%%%%%%%%%%%
% Lab name -- labs are named by the lesson they belong to, never by a
% sequential lab number (Neil, 2026-10-12): this is the L10 Lab.
\def\labLesson{10}

% is this the key or not?
\ifdefined\iskey
	\ifnum\iskey=1
		\keytrue
	\else
		\keyfalse
	\fi
\else
	%% Pick one manually
	\keyfalse % if this is NOT the key
%	\keytrue     % if this IS the key
\fi

% set up the header and footer - change for every term
\ifdefined\thisTermIn
	\def\thisterm{\thisTermIn}
\else
	\def\thisterm{Spring 2027}
\fi
%
\def\thisclass{ECE 448}

%%%%%%%%%%%%%%%%%%%%%%%%%%%%%%%%%%%%%%%%%%%%%%%%%%%%%%
%%% END SET THESE IF VARIABLES
%%%%%%%%%%%%%%%%%%%%%%%%%%%%%%%%%%%%%%%%%%%%%%%%%%%%%%
% packages and shortcuts
\input{myPackages_gr}
\usepackage[hidelinks, linkcolor=red]{hyperref}
\input{myShortcuts}
\setlength{\parindent}{0pt}
\setlength{\parskip}{6pt}

%% exam class customization
\input{exam_config}
% ECE 448 document fonts.
%
% Body text is Barlow (SIL OFL; the faces and the license are in
% latex/fonts/Barlow/). Barlow is not in TeX Live, so it is loaded as an
% OpenType face through fontspec -- which means the document must be built
% with lualatex, not pdflatex. build_practice.sh and build_lab.sh both do.
%
% Compiling with pdflatex anyway is not an error: the guard below falls back
% to the LaTeX default so the source still compiles under any engine and can
% be syntax-checked. The resulting PDF just is not the one we ship.
%
% Paths are relative to latex/, which is where both build scripts run
% (each does cd "$(dirname "$0")"). Loading by path rather than by installed
% family name means no system font install and no fontconfig, so a fresh
% clone builds a byte-comparable PDF anywhere.
%
% MATH IS DELIBERATELY LEFT IN COMPUTER MODERN. Barlow has no math companion,
% and CM math under sans body text is already the house pairing: the decks set
% words in Source Sans Pro and let MathJax render symbols in its CM-derived
% face. Moving math to a sans as well would mean unicode-math plus a sans math
% font (Fira Math being the plausible one) and re-checking every equation in
% the course for symbol coverage. That is a separate decision, not a side
% effect of this one.

\usepackage{iftex}
\ifLuaTeX
	\usepackage{fontspec}
	\defaultfontfeatures{Ligatures=TeX}
	\setmainfont{Barlow}[
		Path           = fonts/Barlow/,
		Extension      = .ttf,
		UprightFont    = *-Regular,
		ItalicFont     = *-Italic,
		BoldFont       = *-Bold,
		BoldItalicFont = *-BoldItalic,
	]
	\setsansfont{Barlow}[
		Path           = fonts/Barlow/,
		Extension      = .ttf,
		UprightFont    = *-Regular,
		ItalicFont     = *-Italic,
		BoldFont       = *-SemiBold,
		BoldItalicFont = *-BoldItalic,
	]
	% microtype earns its place here specifically because of Barlow. Its
	% glyphs are narrower than Computer Modern's, so justified prose sets
	% tighter and TeX has to squeeze interword space to fit; two paragraphs
	% in the L02/L03 sets went overfull on the first Barlow build. Protrusion
	% and font expansion absorb that -- both are lualatex features, so this
	% belongs inside the guard with everything else.
	\usepackage{microtype}

	% Mono stays Latin Modern: Barlow ships no monospaced face, and verbatim
	% and code need one that actually is fixed-pitch.
\fi
   % Barlow body text (needs lualatex)
\printselectedfalse      % full worked solutions in the key, not answers-only
\CorrectChoiceEmphasis{\color{red}}
\SolutionEmphasis{\color{red}}
\renewcommand\questionlabel{\textbf{\thequestion.}}
\renewcommand{\thepartno}{\alph{partno}}
\renewcommand{\partlabel}{(\thepartno)}

% Fillable table cell: blank in the student copy, the expected value in red on
% the KEY. Fixed-width centered columns keep the blank cells writable.
\newcolumntype{K}[1]{>{\centering\arraybackslash}p{#1}}
\newcommand{\kv}[1]{\ifkey{\color{red}#1}\fi\rule[-0.10in]{0pt}{0.34in}}

% Answers are not in this public file. Every solution, \kv value, and \ansbox
% value is \keysol{<id>}, defined in the private ECE448-faculty repo at
% lab-keys/src/ECE448_Lab_L10_PAMLink_solutions.tex; build_lab.sh puts that
% folder on TEXINPUTS (KEYDIR) for the KEY. A KEY build without it stops.
\long\def\defkeysol#1#2{\expandafter\long\expandafter\def\csname keysol@#1\endcsname{#2}}
\def\keysol#1{\ifcsname keysol@#1\endcsname\csname keysol@#1\endcsname\fi}
\ifkey
	\IfFileExists{ECE448_Lab_L10_PAMLink_solutions.tex}%
		{\input{ECE448_Lab_L10_PAMLink_solutions}}%
		{\errmessage{KEY build needs ECE448_Lab_L10_PAMLink_solutions.tex from ECE448-faculty/lab-keys/src (set KEYDIR)}}
\fi

% print the answers if this is the key
\ifkey
	\printanswers % if this IS the key
\else
	\noprintanswers  %if this is NOT the key
\fi

% print the header and footer
% need to print key on top if this is the key
\ifkey
	\firstpageheader{}{\color{red}{ KEY}}{}
	\runningheader{\thisclass\hspace{1pt} -- \thisterm }{L\labLesson\ Lab\hspace{1pt}\color{red}{ KEY}}{\thepage\hspace{1pt} of\hspace{1pt} \numpages}
	\firstpagefooter{}{}{}
	\runningfooter{}{}{}
\else
	\firstpageheader{}{}{}
	\runningheader{\thisclass \hspace{1pt} -- \thisterm }{L\labLesson\ Lab\hspace{1pt}}{\thepage\hspace{1pt} of\hspace{1pt} \numpages}
	\firstpagefooter{}{}{}
	\runningfooter{}{}{}
\fi

\runningheadrule


\begin{document}
\pagestyle{headandfoot}   % exam class

%%%%%%%%%%%%%%%%%%%%%%%%%%%%%%%%%%%%%%%%%%%%%%%%%%%%%%
% title block
%%%%%%%%%%%%%%%%%%%%%%%%%%%%%%%%%%%%%%%%%%%%%%%%%%%%%%

\begin{center}
USAF ACADEMY \\
DEPARTMENT OF ELECTRICAL AND COMPUTER ENGINEERING \\
LESSON \labLesson\ LAB - PAM LINK \\
\textbf{Lab Packet} \\
\thisterm \\[4mm]
\end{center}

\vspace{-4pt}
\textbf{Name:} \rule[-2pt]{2.6in}{0.4pt} \hfill
\textbf{Section:} \rule[-2pt]{0.9in}{0.4pt} \hfill
\textbf{Date:} \rule[-2pt]{1.1in}{0.4pt}

\vspace{4pt}

\fullwidth{\textbf{Learning Objective 4.4: } I can implement a PAM link that transmits and
receives ASCII text end-to-end.}

\vspace{2pt}

This packet is the L\labLesson\ Lab, in three parts, and you turn it in once, complete,
before Lesson 11. Do \textbf{Part 1} on your own before the lab period: run the link
flowgraph and work out the packet you will see at the bench. \textbf{Part 2} is
your own measurements, recorded in class: Questions 3 to 6 on your own laptop, and
Question 7 at a PlutoSDR station with your team. Do \textbf{Part 3} on your own after
the lab: explain, in words, what you saw.

\textbf{Grading.} One submission, graded against Objective 4.4 on the course's
1 / 0 scale, with revision allowed under course policy. Paste screenshots where a
question asks for one.

\textbf{Transmitting.} In class we transmit only into a cable, into a shielded
enclosure, or at very low power in an ISM band, and only when the lesson says to.
Questions 1 to 6 run on your laptop and transmit nothing. Question 7 transmits from a
PlutoSDR's TX port, through a cable with an inline 30 dB SMA attenuator, to the same
Pluto's RX port, at 915 MHz. Connect the cable and attenuator before you run the
flowgraph, and never connect an antenna in this lab. Never set TX attenuation below
30 dB or change the slider's range; the RX input is rated to only a few dBm.

The flowgraphs are on the L10 lesson page, \texttt{L10-pam-tx.grc} and
\texttt{L10-pam-link.grc}, and on the L10 Lab page, \texttt{L10-pam-raw.grc} and
\texttt{L10-pam-pluto.grc}. Every
number below assumes their settings: \texttt{samp\_rate} = 32 kHz, \texttt{sps} = 8,
so 4 kbaud, $\alpha = 0.35$, and the default line \texttt{USAFA ECE 448: ASCII over
PAM.}, 30 characters and a newline, 31 bytes.

\begin{questions}

\fullwidth{\large\textbf{Part 1: Pre-Lab} \normalsize(individual, before the lab period)}

%%%%%%%%%%%%%%%%%%%%%%%%%%%%%%%%%%%%%%%%%%%%%%%%%%%
\question \textbf{Run the link.} Open \texttt{L10-pam-link.grc} in GRC and run it
with the sliders at their defaults, a noise amplitude of 0.3 and a clock error of
500 ppm.

\begin{parts}

	\begin{minipage}[t]{0.58\linewidth}
		\part After about 20 s, record the Lines label.
		\begin{solution}[0.5in]\keysol{sol1}\end{solution}
	\end{minipage}\hfill%
	\begin{minipage}[t]{0.37\linewidth}
		\raggedleft \vspace{12pt}
		\ansbox[1.6in]{\keysol{ans1}}
	\end{minipage}

	\part Paste a screenshot of the running flowgraph's window.
		\begin{solution}[1.6in]\keysol{sol2}\end{solution}

\end{parts}

%%%%%%%%%%%%%%%%%%%%%%%%%%%%%%%%%%%%%%%%%%%%%%%%%%%
\newpage
\question \textbf{The packet.} You will see each of these at the bench.

\begin{parts}

	\begin{minipage}[t]{0.58\linewidth}
		\part Each line is a packet: a 64-bit access code, a 32-bit header, and the 31
		bytes of the line. Find the bits per packet, the time per packet, and the number
		of packets in 30 s.
		\begin{solution}[1.0in]\keysol{sol3}\end{solution}
	\end{minipage}\hfill%
	\begin{minipage}[t]{0.37\linewidth}
		\raggedleft \vspace{12pt}
		bits = \ansbox[1.0in]{\keysol{ans2}}

		\vspace{4pt}
		time = \ansbox[1.0in]{\keysol{ans3}}

		\vspace{4pt}
		in 30 s = \ansbox[1.0in]{\keysol{ans4}}
	\end{minipage}

	\part Choose a line of text of your own, such as your class year and name. Write
	it below, count its bytes with the newline, and write the 16-bit length field the
	header will carry and the 8 bits of its first character.
		\begin{solution}[0.8in]\keysol{sol4}\end{solution}

\end{parts}

%%%%%%%%%%%%%%%%%%%%%%%%%%%%%%%%%%%%%%%%%%%%%%%%%%%
\newpage
\fullwidth{\large\textbf{Part 2: Measurements} \normalsize(in class)}

\question \textbf{Bytes without packets.} Follow the lab page's steps for
\texttt{L10-pam-raw.grc}, the link with no access code and no header. Each run, the
receiver starts listening at a random bit.

\begin{parts}

	\part For three runs, record the bit offset at which the Received bytes label reads
	as text.

	\begin{center}
	\renewcommand{\arraystretch}{1.25}
	\begin{tabular}{| l | K{1.1in} | K{1.1in} | K{1.1in} |}
	\hline
	Run & 1 & 2 & 3 \\\hline
	Offset that reads & \kv{\keysol{kv1}} & \kv{\keysol{kv2}} & \kv{\keysol{kv3}} \\\hline
	\end{tabular}
	\end{center}

		\begin{solution}[0.9in]\keysol{sol5}\end{solution}

	\part Stop and start \texttt{L10-pam-link.grc} three times. Record how long the first
	line takes to appear, and explain in words how the receiver finds the byte boundary
	that you had to find by hand in part (a).
		\begin{solution}[1.3in]\keysol{sol6}\end{solution}

\end{parts}

%%%%%%%%%%%%%%%%%%%%%%%%%%%%%%%%%%%%%%%%%%%%%%%%%%%
\newpage
\question \textbf{Build the transmitter.} Start GRC with an empty canvas. Do not open
\texttt{L10-pam-tx.grc}: build the packet transmitter yourself as the lab page lists,
with your line of text from Question 2(b) in \texttt{msg}.

\begin{parts}

	\part Paste a screenshot of your flowgraph on the GRC canvas. Every connection
	must be black, not red.
		\begin{solution}[1.6in]\keysol{sol7}\end{solution}

	\part Paste a screenshot of the time sink. Mark where the access code ends, where
	the header ends, and the bits of your first character, and check the length field
	against Question 2(b).
		\begin{solution}[1.6in]\keysol{sol8}\end{solution}

	\begin{minipage}[t]{0.58\linewidth}
		\part Read the edges of the shaped signal's spectrum from the frequency sink.
		\begin{solution}[0.6in]\keysol{sol9}\end{solution}
	\end{minipage}\hfill%
	\begin{minipage}[t]{0.37\linewidth}
		\raggedleft \vspace{12pt}
		\ansbox[1.1in]{\keysol{ans5}}
	\end{minipage}

\end{parts}

%%%%%%%%%%%%%%%%%%%%%%%%%%%%%%%%%%%%%%%%%%%%%%%%%%%
\newpage
\question \textbf{Noise.} Run \texttt{L10-pam-link.grc} with the clock error at its
default, 500 ppm, and the matched filter on. Start the noise amplitude at 0.5 and raise
it by 0.1 at a time, waiting about 20 s at each step and watching the Lines label, the
eye, and the Symbol Sync output, which is the 2-PAM constellation drawn against time.
Note the displays at the default 0.3 first. The count restarts each time a slider
moves.

\begin{parts}

	\begin{minipage}[t]{0.58\linewidth}
		\part Record the noise amplitude at which lines start to break, where the correct
		count first falls behind the received count on the Lines label, and the setting
		at which no line arrives intact in 20 s.
		\begin{solution}[1.0in]\keysol{sol10}\end{solution}
	\end{minipage}\hfill%
	\begin{minipage}[t]{0.37\linewidth}
		\raggedleft \vspace{12pt}
		break at \ansbox[1.0in]{\keysol{ans6}}

		\vspace{4pt}
		stop at \ansbox[1.0in]{\keysol{ans7}}
	\end{minipage}

	\part Describe the eye and the Symbol Sync output at 0.3, at your break setting,
	and at your stop setting.

	\begin{center}
	\renewcommand{\arraystretch}{1.25}
	\begin{tabular}{| l | K{1.6in} | K{1.6in} | K{1.6in} |}
	\hline
	Noise & $0.3$ & break & stop \\\hline
	Eye & \kv{\keysol{kv4}} & \kv{\keysol{kv5}} & \kv{\keysol{kv6}} \\\hline
	Symbol Sync output & \kv{\keysol{kv7}} & \kv{\keysol{kv8}} & \kv{\keysol{kv9}} \\\hline
	\end{tabular}
	\end{center}

		\begin{solution}[0.5in]\keysol{sol11}\end{solution}

	\part Near 1.1, wait for a line with a wrong character in the Received label and
	paste a screenshot. Find the character, write its correct and received bytes in
	binary, and count the bits that differ.
		\begin{solution}[1.4in]\keysol{sol12}\end{solution}

	\part Set the noise to 0.5. Switch the matched filter off, wait 20 s, and record
	the Lines label and how the eye looks. Then switch it back on and do the same.

	\begin{center}
	\renewcommand{\arraystretch}{1.25}
	\begin{tabular}{| l | K{2.0in} | K{2.0in} |}
	\hline
	Matched filter & off & on \\\hline
	Lines label & \kv{\keysol{kv10}} & \kv{\keysol{kv11}} \\\hline
	Eye & \kv{\keysol{kv12}} & \kv{\keysol{kv13}} \\\hline
	\end{tabular}
	\end{center}

		\begin{solution}[0.4in]\keysol{sol13}\end{solution}

\end{parts}

%%%%%%%%%%%%%%%%%%%%%%%%%%%%%%%%%%%%%%%%%%%%%%%%%%%
\newpage
\question \textbf{Clock error.} Set the noise amplitude back to 0.3. Set the clock error
to 0, then raise it through the next four values below, in order, without stopping the
flowgraph. At each, wait about 20 s, record the Lines label, and note how the eye sink
looks. Then stop the flowgraph, run it again so the loop locks at the default 500 ppm,
type 40,000 into the clock-error box in one jump, and record the Lines label after
20 s in the last column.

\begin{parts}

	\part Fill in the table.

	\begin{center}
	\renewcommand{\arraystretch}{1.25}
	\begin{tabular}{| l | K{0.55in} | K{0.55in} | K{0.55in} | K{0.55in} | K{0.55in} | K{0.75in} |}
	\hline
	ppm & $0$ & $10{,}000$ & $20{,}000$ & $30{,}000$ & $40{,}000$ & jump to $40{,}000$ \\\hline
	Lines correct & \kv{\keysol{kv14}} & \kv{\keysol{kv15}} & \kv{\keysol{kv16}} & \kv{\keysol{kv17}} & \kv{\keysol{kv18}} & \kv{\keysol{kv19}} \\\hline
	Eye sink & \kv{\keysol{kv20}} & \kv{\keysol{kv21}} & \kv{\keysol{kv22}} & \kv{\keysol{kv23}} & \kv{\keysol{kv24}} & \kv{\keysol{kv25}} \\\hline
	\end{tabular}
	\end{center}

		\begin{solution}[0.9in]\keysol{sol14}\end{solution}

	\part Paste a screenshot at 10,000 ppm, and explain how the eye can look smeared
	while the Symbol Sync output stays at $\pm 1$.
		\begin{solution}[1.2in]\keysol{sol15}\end{solution}

\end{parts}

%%%%%%%%%%%%%%%%%%%%%%%%%%%%%%%%%%%%%%%%%%%%%%%%%%%
\newpage
\question \textbf{The link on a PlutoSDR.} At a PlutoSDR station, with your team, follow
the lab page's steps: connect the Pluto's TX port to its RX port through the 30 dB
attenuator, check
the Pluto with \texttt{iio\_info -s}, and run \texttt{L10-pam-pluto.grc}. The
transmitter and receiver share the Pluto's clock, so there is almost no clock or
frequency error between them; the noise and the carrier phase come from the hardware.

\begin{parts}

	\part If no line arrives within 5 s, switch Polarity. If neither setting works, raise
	the TX attenuation by 10 dB and try both again. Once lines arrive, paste a
	screenshot that shows the constellation, the received line, and the Lines label.
		\begin{solution}[1.6in]\keysol{sol16}\end{solution}

	\part Stop and start the flowgraph three times. Each time, record which polarity
	setting delivers lines.

	\begin{center}
	\renewcommand{\arraystretch}{1.25}
	\begin{tabular}{| l | K{1.1in} | K{1.1in} | K{1.1in} |}
	\hline
	Run & 1 & 2 & 3 \\\hline
	Polarity that works & \kv{\keysol{kv26}} & \kv{\keysol{kv27}} & \kv{\keysol{kv28}} \\\hline
	\end{tabular}
	\end{center}

		\begin{solution}[0.5in]\keysol{sol17}\end{solution}

	\part Stop the flowgraph, put your line from Question 2(b) in \texttt{msg}, run
	it again, and paste a screenshot of your line arriving.
		\begin{solution}[1.2in]\keysol{sol18}\end{solution}

\end{parts}

%%%%%%%%%%%%%%%%%%%%%%%%%%%%%%%%%%%%%%%%%%%%%%%%%%%
\newpage
\fullwidth{\large\textbf{Part 3: Writeup} \normalsize(individual, after the lab)}

\question \textbf{Analysis.} Answer in words; no formulas are needed.

\begin{parts}

	\part Your link went from every line intact to none over a fairly small change in
	the noise amplitude. Explain why, using what a line needs in order to count as
	correct, and say what that means for a link on a real radio.
		\begin{solution}[1.5in]\keysol{sol19}\end{solution}

	\part Explain why switching the matched filter off broke the text at a noise level
	the link handled easily with it on.
		\begin{solution}[1.5in]\keysol{sol20}\end{solution}

	\part Your link held nearly every line to at least 10,000 ppm. Compare that with two
	radios whose crystals are accurate to 20 ppm each, and say what would happen to the
	link if you replaced Symbol Sync with Keep 1 in N, which samples at a fixed point.
		\begin{solution}[1.5in]\keysol{sol21}\end{solution}

	\part Explain why the working polarity changed between runs on the Pluto, and
	propose a change to the receiver that would make the polarity switch unnecessary.
		\begin{solution}[1.5in]\keysol{sol22}\end{solution}

\end{parts}

\end{questions}

\vspace{1.0em}
\noindent\textbf{Documentation:} \rule[-2pt]{0.6\linewidth}{0.4pt}

\end{document}
